% T12最小订正片段；供择段粘贴。需要amsmath；文献键由整合者按实际bib补入。
% 依据：Duguet et al. (2015), PDF pp. 4--8, Eqs. (4)--(5), (9)--(15), (17)--(18), (23)--(29).
令 $U_\mu^p=\langle\Psi_0^A|c_p|\Psi_\mu^{A+1}\rangle$、
$V_\nu^p=\langle\Psi_0^A|c_p^\dagger|\Psi_\nu^{A-1}\rangle$，则
\begin{align}
(S_\mu^+)_{pq}&=U_\mu^p(U_\mu^q)^*, &
(S_\nu^-)_{pq}&=(V_\nu^p)^*V_\nu^q,\\
h^{\rm cent}_{pq}&=\langle\Psi_0^A|\{[c_p,H],c_q^\dagger\}|\Psi_0^A\rangle\\
&=\sum_\mu(S_\mu^+)_{pq}E_\mu^+
 +\sum_\nu(S_\nu^-)_{pq}E_\nu^- .
\end{align}
Baranger ESPE 是重心矩阵 $h^{\rm cent}$ 的本征值。在球形 $0^+$ 参考态下，矩阵按 $j,m$、宇称和粒子种类分块；同一块内不同径向轨道仍可混合，须先对角化，才能把本征值写成该本征基底中添、去粒谱强度的加权和。

在固定基底、球对称密度与两体相互作用的特定模型中，若将 $V_{jj}$ 定义为同种粒子 $j$ 壳内 $g_j(g_j-1)$ 个反对称有序磁子态对的平均，$g_j=2j+1>1$，则同轨密度贡献为 $V_{jj}(1-1/g_j)n_j$；相对 $V_{jj}n_j$ 的差为 $-V_{jj}n_j/g_j$。将其直接用于某个 ESPE 本征值还须保证所选轨道基底使重心矩阵对角，并另计三体力等贡献。对于相同相互作用和基底的两态，该差的变化为 $-V_{jj}\Delta n_j/g_j$；是否可忽略须结合占据数变化和相互作用强度判断。

若 SRG 演化哈密顿量与态而保留裸单核子场算符的形式，谱因子和重心通常具有分辨率及方案依赖；对场算符也作一致的精确幺正变换，则相应振幅与重心保持不变，但演化后的场一般包含高体项。具体流或个别能级可以不变。实验能谱、分离能和截面可测；由实验数据按明确规则构造的加权重心亦可报告，其能窗、权重及不确定度应写明；由反应模型提取的谱因子加权重心还依赖该模型，不能不加限定地等同于完整 Baranger ESPE。

比较不同计算的 ESPE 时，应交代哈密顿量、分辨率与方案、场算符、单粒子基底、多体截断、谱强度覆盖和不确定度。$\lambda=1.8$--$2.2\,\mathrm{fm}^{-1}$、NO2B 或动力学自能平缓可作为具体计算的选择或待检验近似，而非一般比较的必要充分条件。

% 载体注：2026-11-08单次文章PDF请求返回未获entitlement且限第一页；
% 不可据此推断2026-11-13权限或整个过刊库。现存Baranger PDF仅1页；XML rawtext是有损OCR。
